\Chapter{29}

\Verse{1}
{
    \zA{话说宝玉正自发怔，不想黛玉将手帕子扔了来，正碰在眼睛上，倒唬了一跳， 问：“这是谁？” 黛玉摇着头儿笑道：“不敢，是我失了手。 因为宝姐姐要看呆雁，
        我比给他看，不想失了手。” 宝玉揉着眼睛，待要说什么，又不好说的。 一时凤姐儿来了。 因说起初一日在清虚观打醮的事来，约着宝钗、宝玉、黛玉 等看戏去。
        宝钗笑道：“罢，罢，怪热的，什么没看过的戏!我不去。” 凤姐道： “他们那里凉快，两边又有楼。}
}
{
    \eA{Pao-yü, so our story runs, was gazing vacantly, when Tai-yü, at a moment
        least expected, flung her handkerchief at him, which just hit him on the
        eyes, and frightened him out of his wits. "Who was it?" he cried. Lin Tai-yü
        nodded her head and smiled. "I would not venture to do such a thing," she
        said, "it was a mere slip of my hand.}
}
{
}

\Verse{2}
{
    \zA{咱们要去，我头几天先打发人去，把那些道士都赶 出去，把楼上打扫了，挂起帘子来，一个闲人不许放进庙去，才是好呢。 我已经回 了太太了，你们不去，我自家去。
        这些日子也闷的很了，家里唱动戏，我又不得舒 舒服服的看。” 贾母听说，就笑道：“既这么着，我和你去。” 凤姐听说，笑道： “老祖宗也去？
        敢仔好，可就是我又不得受用了。” 贾母道：“到明儿我在正面楼 上，你在傍边楼上，你也不用到我这边来立规矩，可好不好？”}
}
{
    \eA{As cousin Pao-ch'ai wished to see the silly wild goose, I was pointing it
        out to her, when the handkerchief inadvertently flew out of my grip." Pao-yü
        kept on rubbing his eyes. The idea suggested itself to him to make some
        remonstrance, but he could not again very well open his lips. Presently,
        lady Feng arrived.}
}
{
}

\Verse{3}
{
    \zA{凤姐笑道：“这就 是老祖宗疼我了。” 贾母因向宝钗道：“你也去，连你母亲也去； 长天老日的，在 家里也是睡觉。” 宝钗只得答应着。
        贾母又打发人去请了薛姨妈，顺路告诉王夫人，要带了他们姊妹去。 王夫人因 一则身上不好，二则预备元春有人出来，早已回了不去的； 听贾母如此说，笑道：
        “还是这么高兴。 打发人去到园里告诉，有要逛去的，只管初一跟老太太逛去。”}
}
{
    \eA{She then alluded, in the course of conversation, to the thanksgiving
        service, which was to be offered on the first, in the Ch'ing Hsü temple, and
        invited Pao-ch'ai, Pao-yü, Tai-yü and the other inmates with them to be
        present at the theatricals. "Never mind," smiled Pao-ch'ai, "it's too hot;
        besides, what plays haven't I seen? I don't mean to come." "It's cool enough
        over at their place," answered lady Feng.}
}
{
}

\Verse{4}
{
    \zA{这个话一传开了，别人还可已，只是那些丫头们，天天不得出门槛儿，听了这话谁 不要去，就是各人的主子懒怠去，他也百般的撺掇了去：因此李纨等都说去。 贾母
        心中越发喜欢，早已吩咐人去打扫安置，不必细说。 单表到了初一这一日，荣国府门前车辆纷纷，人马簇簇，那底下执事人等，听
        见是贵妃做好事，贾母亲去拈香，况是端阳佳节，因此凡动用的物件，一色都是齐 全的，不同往日。}
}
{
    \eA{"There are also two-storied buildings on either side; so we must all go!
        I'll send servants a few days before to drive all that herd of Taoist
        priests out, to sweep the upper stories, hang up curtains, and to keep out
        every single loafer from the interior of the temple; so it will be all right
        like that.}
}
{
}

\Verse{5}
{
    \zA{少时贾母等出来，贾母坐一乘八人大轿，李氏、凤姐、薛姨妈每 人一乘四人轿，宝钗、黛玉二人共坐一辆翠盖珠缨八宝车，迎春、探春、惜春三人 共坐一辆朱轮华盖车。}
}
{
    \eA{I've already told our Madame Wang that if you people don't go, I mean to go
        all alone, as I've been again in very low spirits these last few days, and
        as when theatricals come off at home, it's out of the question for me to
        look on with any peace and quiet." When dowager lady Chia heard what she
        said, she smiled. "Well, in that case," she remarked, "I'll go along with
        you." Lady Feng, at these words, gave a smile.}
}
{
}

\Verse{6}
{
    \zA{然后贾母的丫头鸳鸯、鹦鹉、琥珀、珍珠，黛玉的丫头紫鹃、 雪雁、鹦哥，宝钗的丫头莺儿、文杏，迎春的丫头司棋、绣橘，探春的丫头侍书、
        翠墨，惜春的丫头入画、彩屏，薛姨妈的丫头同喜、同贵，外带香菱，香菱的丫头 臻儿，李氏的丫头素云、碧月，凤姐儿的丫头平儿、丰儿、小红，并王夫人的两个
        丫头金钏、彩云，也跟了凤姐儿来。 奶子抱着大姐儿，另在一辆车上。 还有几个粗 使的丫头，连上各房的老嬷嬷奶妈子，并跟着出门的媳妇子们，黑压压的站了一街
        的车。}
}
{
    \eA{"Venerable ancestor," she replied, "were you also to go, it would be ever so
        much better; yet I won't feel quite at my ease!" "To-morrow," dowager lady
        Chia continued, "I can stay in the two-storied building, situated on the
        principal site, while you can go to the one on the side. You can then
        likewise dispense with coming over to where I shall be to stand on any
        ceremonies. Will this suit you or not?"}
}
{
}

\Verse{7}
{
    \zA{那街上的人见是贾府去烧香，都站在两边观看。 那些小门小户的妇女，也都 开了门在门口站着，七言八语，指手画脚，就像看那过会的一般。 只见前头的全副
        执事摆开，一位青年公子骑着银鞍白马，彩辔朱缨，在那八人轿前领着那些车轿人 马，浩浩荡荡，一片锦绣香烟，遮天压地而来。 却是鸦雀无闻，只有车轮马蹄之声。
        不多时，已到了清虚观门口。 只听钟鸣鼓响，早有张法官执香披衣，带领众道 士在路旁迎接。}
}
{
    \eA{"This is indeed," lady Feng smiled, "a proof of your regard for me, my
        worthy senior." Old lady Chia at this stage faced Pao-ch'ai. "You too should
        go," she said, "so should your mother; for if you remain the whole day long
        at home, you will again sleep your head off." Pao-ch'ai felt constrained to
        signify her assent. Dowager lady Chia then also despatched domestics to
        invite Mrs. Hsüeh;}
}
{
}

\Verse{8}
{
    \zA{宝玉下了马，贾母的轿刚至山门以内，见了本境城隍土地各位泥塑 圣像，便命住轿。 贾珍带领各子弟上来迎接。 凤姐儿的轿子却赶在头里先到了，带
        着鸳鸯等迎接上来，见贾母下了轿，忙要搀扶。 可巧有个十二三岁的小道士儿，拿 着个剪筒，照管各处剪蜡花儿，正欲得便且藏出去，不想一头撞在凤姐儿怀里。 凤
        姐便一扬手照脸打了个嘴巴，把那小孩子打了一个斤斗，骂道：“小野杂种!往那 里跑？” 那小道士也不顾拾烛剪，爬起来往外还要跑。}
}
{
    \eA{and, on their way, they notified Madame Wang that she was to take the young
        ladies along with her. But Madame Wang felt, in the first place, in a poor
        state of health, and was, in the second, engaged in making preparations for
        the reception of any arrivals from Yüan Ch'un, so that she, at an early
        hour, sent word that it was impossible for her to leave the house.}
}
{
}

\Verse{9}
{
    \zA{正值宝钗等下车，众婆娘媳 妇正围随的风雨不透，但见一个小道士滚了出来，都喝声叫：“拿，拿!打，打！” 贾母听了，忙问：“是怎么了？” 贾珍忙过来问。
        凤姐上去搀住贾母，就回说：“一 个小道士儿剪蜡花的，没躲出去，这会子混钻呢。” 贾母听说，忙道：“快带了那 孩子来，别唬着他。
        小门小户的孩子，都是娇生惯养惯了的，那里见过这个势派？ 倘或唬着他，倒怪可怜见儿的。 他老子娘岂不疼呢。” 说着，便叫贾珍去好生带了 来。}
}
{
    \eA{Yet when she received old lady Chia's behest, she smiled and exclaimed: "Are
        her spirits still so buoyant!" and transmitted the message into the garden
        that any, who had any wish to avail themselves of the opportunity, were at
        liberty to go on the first, with their venerable senior as their chaperonne.
        As soon as these tidings were spread abroad, every one else was indifferent
        as to whether they went or not;}
}
{
}

\Verse{10}
{
    \zA{贾珍只得去拉了，那孩子一手拿着蜡剪，跪在地下乱颤。 贾母命贾珍拉起来， 叫他不用怕，问他几岁了。 那孩子总说不出话来。 贾母还说：“可怜见儿的！” 又
        向贾珍道：“珍哥带他去罢。 给他几个钱买果子吃，别叫人难为了他。” 贾珍答应， 领出去了。 这里贾母带着众人，一层一层的瞻拜观玩。
        外面小厮们见贾母等进入二层山 门，忽见贾珍领了个小道士出来，叫人：“来带了去，给他几百钱，别难为了他。” 家人听说，忙上来领去。}
}
{
    \eA{but of those girls who, day after day, never put their foot outside the
        doorstep, which of them was not keen upon going, the moment they heard the
        permission conceded to them? Even if any of their respective mistresses were
        too lazy to move, they employed every expedient to induce them to go. Hence
        it was that Li Kung-ts'ai and the other inmates signified their unanimous
        intention to be present.}
}
{
}

\Verse{11}
{
    \zA{贾珍站在台阶上，因问：“管家在那里？” 底下站的小厮 们见问，都一齐喝声说：“叫管家！” 登时林之孝一手整理着帽子，跑进来，到了 贾珍跟前。
        贾珍道：“虽然这里地方儿大，今儿咱们的人多，你使的人，你就带了 在这院里罢，使不着的，打发到那院里去。 把小么儿们多挑几个在这二层门上和两
        边的角门上，伺候着要东西传话。 你可知道不知道？ 今儿姑娘奶奶们都出来，一个 闲人也不许到这里来。” 林之孝忙答应“知道”，又说了几个“是”。}
}
{
    \eA{Dowager lady Chia, at this, grew more exultant than ever, and she issued
        immediate directions for servants to go and sweep and put things in proper
        order. But to all these preparations, there is no necessity of making
        detailed reference; sufficient to relate that on the first day of the moon,
        carriages stood in a thick maze, and men and horses in close concourse, at
        the entrance of the Jung Kuo mansion.}
}
{
}

\Verse{12}
{
    \zA{贾珍道：“去 罢。” 又问：“怎么不见蓉儿？” 一声未了，只见贾蓉从钟楼里跑出来了。 贾珍道： “你瞧瞧，我这里没热，他倒凉快去了！” 喝命家人啐他。
        那小厮们都知道贾珍素 日的性子，违拗不得，就有个小厮上来向贾蓉脸上啐了一口。 贾珍还瞪着他，那小 厮便问贾蓉：“爷还不怕热，哥儿怎么先凉快去了？”
        贾蓉垂着手，一声不敢言语。 那贾芸、贾萍、贾芹等听见了，不但他们慌了，并贾琏、贾？}
}
{
    \eA{When the servants, the various managers and other domestics came to learn
        that the Imperial Consort was to perform good deeds and that dowager lady
        Chia was to go in person and offer incense, they arranged, as it happened
        that the first of the moon, which was the principal day of the ceremonies,
        was, in addition, the season of the dragon-boat festival, all the necessary
        articles in perfect readiness and with unusual splendour.}
}
{
}

\Verse{13}
{
    \zA{、贾琼等也都忙了， 一个一个都从墙根儿底下慢慢的溜下来了。 贾珍又向贾蓉道：“你站着做什么？ 还
        不骑了马跑到家里告诉你娘母子去!老太太和姑娘们都来了，叫他们快来伺候！” 贾蓉听说，忙跑了出来，一叠连声的要马。 一面抱怨道：“早都不知做什么的，这
        会子寻趁我。” 一面又骂小子：“捆着手呢么？ 马也拉不来！” 要打发小厮去，又 恐怕后来对出来，说不得亲自走一趟，骑马去了。}
}
{
    \eA{Shortly, old lady Chia and the other inmates started on their way. The old
        lady sat in an official chair, carried by eight bearers: widow Li, lady Feng
        and Mrs. Hsüeh, each in a four-bearer chair. Pao-ch'ai and Tai-yü mounted
        together a curricle with green cover and pearl tassels, bearing the eight
        precious things.}
}
{
}

\Verse{14}
{
    \zA{且说贾珍方要抽身进来，只见张道士站在傍边，陪笑说道：“论理，我不比别 人，应该里头伺候； 只因天气炎热，众位千金都出来了，法官不敢擅入，请爷的示 下。
        恐老太太问，或要随喜那里，我只在这里伺候罢了。” 贾珍知道这张道士虽然 是当日荣国公的替身，曾经先皇御口亲呼为“大幻仙人”，如今现掌道录司印，又
        是当今封为“终了真人”，现今王公藩镇都称为神仙，所以不敢轻慢。 二则他又常 往两个府里去，太太姑娘们都是见的。}
}
{
    \eA{The three sisters, Ying Ch'un, T'an Ch'un, and Hsi Ch'un got in a carriage
        with red wheels and ornamented hood. Next in order, followed dowager lady
        Chia's waiting-maids, Yüan Yang, Ying Wu, Hu Po, Chen Chu; Lin Tai-yü's
        waiting-maids Tzu Chüan, Hsüeh Yen, and Ch'un Ch'ien; Pao-ch'ai's
        waiting-maids Ying Erh and Wen Hsing; Ying Ch'un's servant-girls Ssu Ch'i
        and Hsiu Chü;}
}
{
}

\Verse{15}
{
    \zA{今见他如此说，便笑道：“咱们自己，你又 说起这话来。 再多说，我把你这胡子还揪了你的呢!还不跟我进来呢。” 那张道士 呵呵的笑着，跟了贾珍进来。
        贾珍到贾母跟前，控身陪笑，说道：“张爷爷进来请安。” 贾母听了，忙道： “请他来。” 贾珍忙去搀过来。 那张道士先呵呵笑道：“无量寿佛!老祖宗一向福
        寿康宁，众位奶奶姑娘纳福!一向没到府里请安，老太太气色越发好了。” 贾母笑 道：“老神仙你好？” 张道士笑道：“托老太太的万福，小道也还康健。}
}
{
    \eA{T'an Ch'un's waiting-maids Shih Shu and Ts'ui Mo; Hsi Ch'un's servant-girls
        Ju Hua and Ts'ai P'ing; and Mrs. Hsüeh's waiting-maids T'ung Hsi, and T'ung
        Kuei. Besides these, were joined to their retinue: Hsiang Ling and Hsiang
        Ling's servant-girl Ch'in Erh; Mrs. Li's waiting-maids Su Yün and Pi Yüeh;}
}
{
}

\Verse{16}
{
    \zA{别的倒罢 了，只记挂着哥儿，一向身上好？ 前日四月二十六，我这里做遮天大王的圣诞，人 也来的少，东西也很干净，我说请哥儿来逛逛，怎么说不在家？”
        贾母说道：“果 真不在家。” 一面回头叫宝玉。 谁知宝玉解手儿去了，才来，忙上前问：“张爷爷好？” 张道士也抱住问了好，
        又向贾母笑道：“哥儿越发发福了。” 贾母道：“他外头好，里头弱。 又搭着他老 子逼着他念书，生生儿的把个孩子逼出病来了。”}
}
{
    \eA{lady Feng's servant-girls P'ing Erh, Feng Erh and Hsiao Hung, as well as
        Madame Wang's two waiting-maids Chin Ch'uan and Ts'ai Yün. Along with lady
        Feng, came a nurse carrying Ta Chieh Erh. She drove in a separate carriage,
        together with a couple of servant-girls.}
}
{
}

\Verse{17}
{
    \zA{张道士道：“前日我在好几处看 见哥儿写的字，做的诗，都好的了不得。 怎么老爷还抱怨哥儿不大喜欢念书呢？ 依 小道看来，也就罢了。”
        又叹道：“我看见哥儿的这个形容身段，言谈举动，怎么 就和当日国公爷一个稿子！” 说着，两眼酸酸的。 贾母听了，也由不得有些戚惨， 说道：“正是呢。
        我养了这些儿子孙子，也没一个像他爷爷的，就只这玉儿还像他 爷爷。” 那张道士又向贾珍道：“当日国公爷的模样儿，爷们一辈儿的不用说了， 自然没赶上；}
}
{
    \eA{Added also to the number of the suite were matrons and nurses, attached to
        the various establishments, and the wives of the servants of the household,
        who were in attendance out of doors. Their carriages, forming one black
        solid mass, therefore, crammed the whole extent of the street.}
}
{
}

\Verse{18}
{
    \zA{大约连大老爷、二老爷也记不清楚了罢？” 说毕，又呵呵大笑道：“前 日在一个人家儿，看见位小姐，今年十五岁了，长的倒也好个模样儿。 我想着哥儿
        也该提亲了。 要论这小姐的模样儿，聪明智慧，根基家当，倒也配的过。 但不知老 太太怎么样？ 小道也不敢造次，等请了示下，才敢提去呢。”
        贾母道：“上回有个 和尚说了，这孩子命里不该早娶，等再大一大儿再定罢。 你如今也讯听着，不管他 根基富贵，只要模样儿配的上，就来告诉我。}
}
{
    \eA{Dowager lady Chia and other members of the party had already proceeded a
        considerable distance in their chairs, and yet the inmates at the gate had
        not finished mounting their vehicles. This one shouted: "I won't sit with
        you." That one cried: "You've crushed our mistress' bundle." In the
        carriages yonder, one screamed: "You've pulled my flowers off." Another one
        nearer exclaimed: "You've broken my fan."}
}
{
}

\Verse{19}
{
    \zA{就是那家子穷，也不过帮他几两银子 就完了。 只是模样儿性格儿难得好的。” 说毕，只见凤姐儿笑道：“张爷爷，我们丫头的寄名符儿，你也不换去，前儿
        亏你还有那么大脸，打发人和我要鹅黄缎子去!要不给你，又恐怕你那老脸上下不 来。” 张道士哈哈大笑道：“你瞧，我眼花了!也没见奶奶在这里，也没道谢。 寄
        名符早已有了，前日原想送去，不承望娘娘来做好事，也就混忘了。 还在佛前镇着 呢。 等着我取了来。”}
}
{
    \eA{And they chatted and chatted, and talked and laughed with such incessant
        volubility, that Chou Jui's wife had to go backward and forward calling them
        to task. "Girls," she said, "this is the street. The on-lookers will laugh
        at you!" But it was only after she had expostulated with them several times
        that any sign of improvement became at last visible.}
}
{
}

\Verse{20}
{
    \zA{说着，跑到大殿上，一时拿了个茶盘，搭着大红蟒缎经袱子， 托出符来。 大姐儿的奶子接了符。 张道士才要抱过大姐儿来，只见凤姐笑道：“你
        就手里拿出来罢了，又拿个盘子托着！” 张道士道：“手里不干不净的，怎么拿？ 用盘子洁净些。” 凤姐笑道：“你只顾拿出盘子，倒唬了我一跳。 我不说你是为送
        符，倒像和我们化布施来了。” 众人听说哄然一笑，连贾珍也掌不住笑了。 贾母回 头道：“猴儿，猴儿!你不怕下割舌地狱？” 凤姐笑道：“我们爷儿们不相干。}
}
{
    \eA{The van of the procession had long ago reached the entrance of the Ch'ing
        Hsü Temple. Pao-yü rode on horseback. He preceded the chair occupied by his
        grandmother Chia. The throngs that filled the streets ranged themselves on
        either side. On their arrival at the temple, the sound of bells and the
        rattle of drums struck their ear. Forthwith appeared the head-bonze Chang, a
        stick of incense in hand;}
}
{
}

\Verse{21}
{
    \zA{他 怎么常常的说我该积阴骘、迟了就短命呢？” 张道士也笑道：“我拿出盘子来，一 举两用，倒不为化布施，倒要把哥儿的那块玉请下来，托出去给那些远来的道友和
        徒子徒孙们见识见识。” 贾母道：“既这么着，你老人家老天拔地的，跑什么呢， 带着他去瞧了叫他进来，就是了。” 张道士道：“老太太不知道，看着小道是八十
        岁的人，托老太太的福，倒还硬朗； 二则外头的人多气味难闻，况且大暑热的天， 哥儿受不惯，倘或哥儿中了腌？ 气味，倒值多了。”}
}
{
    \eA{his cloak thrown over his shoulders. He took his stand by the wayside at the
        head of a company of Taoist priests to present his greetings. The moment
        dowager lady Chia reached, in her chair, the interior of the main gate, she
        descried the lares and penates, the lord presiding over that particular
        district, and the clay images of the various gods, and she at once gave
        orders to halt.}
}
{
}

\Verse{22}
{
    \zA{贾母听说，便命宝玉摘下通灵 玉来，放在盘内。 那张道士兢兢业业的用蟒袱子垫着，捧出去了。 这里贾母带着众人各处游玩一回，方去上楼。
        只见贾珍回说：“张爷爷送了玉 来。” 刚说着，张道士捧着盘子走到跟前，笑道：“众人托小道的福，见了哥儿的
        玉，实在稀罕，都没什么敬贺的，这是他们各人传道的法器，都愿意为敬贺之礼。 虽不稀罕，哥儿只留着玩耍赏人罢。” 贾母听说，向盘内看时，只见也有金璜，也
        有玉？}
}
{
    \eA{Chia Chen advanced to receive her acting as leader to the male members of
        the family. Lady Feng was well aware that Yüan Yang and the other attendants
        were at the back and could not overtake their old mistress, so she herself
        alighted from her chair to volunteer her services.}
}
{
}

\Verse{23}
{
    \zA{，或有“事事如意”，或有“岁岁平安”，皆是珠穿宝嵌、玉琢金镂，共有 三五十件。 因说道：“你也胡闹。 他们出家人，是那里来的？ 何必这样？
        这断不能收。” 张道士笑道：“这是他们一点敬意，小道也不能阻挡。 老太太要不留下，倒叫他们 看着小道微薄，不像是门下出身了。”
        贾母听如此说，方命人接下了。 宝玉笑道： “老太太，张爷爷既这么说，又推辞不得，我要这个也无用，不如叫小子捧了这个， 跟着我出去散给穷人罢。”
        贾母笑道：“这话说的也是。”}
}
{
    \eA{She was about to hastily press forward and support her, when, by a strange
        accident, a young Taoist neophyte, of twelve or thirteen years of age, who
        held a case containing scissors, with which he had been snuffing the candles
        burning in the various places, just seized the opportunity to run out and
        hide himself, when he unawares rushed, head foremost, into lady Feng's arms.}
}
{
}

\Verse{24}
{
    \zA{张道士忙拦道：“哥儿 虽要行好，但这些东西虽说不甚稀罕，也到底是几件器皿。 若给了穷人，一则与他 们也无益，二则反倒遭塌了这些东西。
        要舍给穷人，何不就散钱给他们呢？” 宝玉 听说，便命：“收下，等晚上拿钱施舍罢。” 说毕，张道士方才退出。 这里贾母和众人上了楼，在正面楼上归坐。
        凤姐等上了东楼。 众丫头等在西楼 轮流伺候。 一时贾珍上来回道：“神前拈了戏，头一本是《白蛇记》。” 贾母便问： “是什么故事？”}
}
{
    \eA{Lady Feng speedily raised her hand and gave him such a slap on the face that
        she made the young fellow reel over and perform a somersault. "You boorish
        young bastard!" she shouted, "where are you running to?" The young Taoist
        did not even give a thought to picking up the scissors, but crawling up on
        to his feet again, he tried to scamper outside.}
}
{
}

\Verse{25}
{
    \zA{贾珍道：“汉高祖斩蛇起首的故事。 第二本是《满床笏》。” 贾 母点头道：“倒是第二本？ 也还罢了。 神佛既这样，也只得如此。” 又问：“第三 本？”
        贾珍道：“第三本是《南柯梦》。” 贾母听了，便不言语。 贾珍退下来，走 至外边，预备着申表、焚钱粮、开戏，不在话下。
        且说宝玉在楼上，坐在贾母傍边，因叫个小丫头子捧着方才那一盘子东西，将 自己的玉带上，用手翻弄寻拨，一件一件的挑与贾母看。}
}
{
    \eA{But just at that very moment Pao-ch'ai and the rest of the young ladies were
        dismounting from their vehicles, and the matrons and women-servants were
        closing them in so thoroughly on all sides that not a puff of wind or a drop
        of rain could penetrate, and when they perceived a Taoist neophyte come
        rushing headlong out of the place, they, with one voice, exclaimed: "Catch
        him, catch him! Beat him, beat him!"}
}
{
}

\Verse{26}
{
    \zA{贾母因看见有个赤金点翠 的麒麟，便伸手拿起来，笑道：“这件东西，好像是我看见谁家的孩子也带着一个 的。” 宝钗笑道：“史大妹妹有一个，比这个小些。”
        贾母道：“原来是云儿有这 个。” 宝玉道：“他这么往我们家去住着，我也没看见？” 探春笑道：“宝姐姐有 心，不管什么他都记得。”
        黛玉冷笑道：“他在别的上头心还有限，惟有这些人带 的东西上，他才是留心呢。” 宝钗听说，回头装没听见。}
}
{
    \eA{Old lady Chia overheard their cries. She asked with alacrity what the fuss
        was all about. Chia Chen immediately stepped outside to make inquiries. Lady
        Feng then advanced and, propping up her old senior, she went on to explain
        to her that a young Taoist priest, whose duties were to snuff the candles,
        had not previously retired out of the compound, and that he was now
        endeavouring to recklessly force his way out."}
}
{
}

\Verse{27}
{
    \zA{宝玉听见史湘云有这件东 西，自己便将那麒麟忙拿起来，揣在怀里。 忽又想到怕人看见他听是史湘云有了， 他就留着这件，因此手里揣着，却拿眼睛瞟人。
        只见众人倒都不理论，惟有黛玉瞅 着他点头儿，似有赞叹之意。 宝玉心里不觉没意思起来，又掏出来，瞅着黛玉讪笑
        道：“这个东西有趣儿，我替你拿着，到家里穿上个穗子你带，好不好？” 黛玉将 头一扭道：“我不稀罕。” 宝玉笑道：“你既不稀罕，我可就拿着了。” 说着，又
        揣起来。}
}
{
    \eA{"Be quick and bring the lad here," shouted dowager lady Chia, as soon as she
        heard her explanation, "but, mind, don't frighten him. Children of mean
        families invariably get into the way of being spoilt by over-indulgence. How
        ever could he have set eyes before upon such display as this! Were you to
        frighten him, he will really be much to be pitied; and won't his father and
        mother be exceedingly cut up?"}
}
{
}

\Verse{28}
{
    \zA{刚要说话，只见贾珍之妻尤氏和贾蓉续娶的媳妇胡氏，婆媳两个来了，见过贾 母。 贾母道：“你们又来做什么，我不过没事来逛逛。” 一句话说了，只见人报：
        “冯将军家有人来了。” 原来冯紫英家听见贾府在庙里打醮，连忙预备猪羊、香烛、 茶食之类，赶来送礼。 凤姐听了，忙赶过正楼来，拍手笑道：“嗳呀!我却没防着
        这个。 只说咱们娘儿们来闲逛逛，人家只当咱们大摆斋坛的来送礼。 都是老太太闹 的!这又不得预备赏封儿。”}
}
{
    \eA{As she spoke, she asked Chia Chen to go and do his best to bring him round.
        Chia Chen felt under the necessity of going, and he managed to drag the lad
        into her presence. With the scissors still clasped in his hand, the lad fell
        on his knees, and trembled violently. Dowager lady Chia bade Chia Chen raise
        him up. "There's nothing to fear!" she said reassuringly. Then she asked him
        how old he was.}
}
{
}

\Verse{29}
{
    \zA{刚说了，只见冯家的两个管家女人上楼来了。 冯家两 个未去，接着赵侍郎家也有礼来了。 于是接二连三，都听见贾府打醮，女眷都在庙
        里，凡一应远亲近友，世家相与，都来送礼。 贾母才后悔起来，说：“又不是什么 正经斋事，我们不过闲逛逛，没的惊动人。” 因此虽看了一天戏，至下午便回来了。
        次日便懒怠去。 凤姐又说：“‘打墙也是动土’，已经惊动了人，今儿乐得还去逛 逛。”}
}
{
    \eA{The boy, however, could on no account give vent to speech. "Poor boy!" once
        more exclaimed the old lady. And continuing: "Brother Chen," she added,
        addressing herself to Chia Chen, "take him away, and give him a few cash to
        buy himself fruit with; and do impress upon every one that they are not to
        bully him." Chia Chen signified his assent and led him off.}
}
{
}

\Verse{30}
{
    \zA{贾母因昨日见张道士提起宝玉说亲的事来，谁知宝玉一日心中不自在，回家 来生气，嗔着张道士与他说了亲，口口声声说“从今以后，再不见张道士了”，别
        人也并不知为什么原故。 二则黛玉昨日回家，又中了暑。 因此二事，贾母便执意不 去了。 凤姐见不去，自己带了人去，也不在话下。
        且说宝玉因见黛玉病了，心里放不下，饭也懒怠吃，不时来问，只怕他有个好 歹。 黛玉因说道：“你只管听你的戏去罢，在家里做什么？”}
}
{
    \eA{During this time, old lady Chia, taking along with her the whole family
        party, paid her devotions in storey after storey, and visited every place.
        The young pages, who stood outside, watched their old mistress and the other
        inmates enter the second row of gates. But of a sudden they espied Chia Chen
        wend his way outwards, leading a young Taoist priest, and calling the
        servants to come, say;}
}
{
}

\Verse{31}
{
    \zA{宝玉因昨日张道士提 亲之事，心中大不受用，今听见黛玉如此说，心里因想道：“别人不知道我的心还 可恕，连他也奚落起我来。”
        因此心中更比往日的烦恼加了百倍。 要是别人跟前， 断不能动这肝火，只是黛玉说了这话，倒又比往日别人说这话不同，由不得立刻沉
        下脸来，说道：“我白认得你了!罢了，罢了！” 黛玉听说，冷笑了两声道：“你 白认得了我吗？ 我那里能够像人家有什么配的上你的呢！”}
}
{
    \eA{"Take him and give him several hundreds of cash and abstain from
        ill-treating him." At these orders, the domestics approached with hurried
        step and led him off. Chia Chen then inquired from the terrace-steps where
        the majordomo was. At this inquiry, the pages standing below, called out in
        chorus, "Majordomo!"}
}
{
}

\Verse{32}
{
    \zA{宝玉听了，便走来，直 问到脸上道：“你这么说，是安心咒我天诛地灭？” 黛玉一时解不过这话来。 宝玉
        又道：“昨儿还为这个起了誓呢，今儿你到底儿又重我一句!我就天诛地灭，你又 有什么益处呢？” 黛玉一闻此言，方想起昨日的话来。 今日原自己说错了，又是急，
        又是愧，便抽抽搭搭的哭起来，说道：“我要安心咒你，我也天诛地灭!何苦来呢! 我知道昨日张道士说亲，你怕拦了你的好姻缘，你心里生气，来拿我煞性子！”}
}
{
    \eA{Lin Chih-hsiao ran over at once, while adjusting his hat with one hand, and
        appeared in the presence of Chia Chen. "Albeit this is a spacious place,"
        Chia Chen began, "we muster a good concourse to-day, so you'd better bring
        into this court those servants, who'll be of any use to you, and send over
        into that one those who won't.}
}
{
}

\Verse{33}
{
    \zA{原来宝玉自幼生成来的有一种下流痴病，况从幼时和黛玉耳鬓厮磨，心情相 对，如今稍知些事，又看了些邪书僻传，凡远亲近友之家所见的那些闺英闱秀，皆
        未有稍及黛玉者，所以早存一段心事，只不好说出来。 故每每或喜或怒，变尽法子 暗中试探。 那黛玉偏生也是个有些痴病的，也每用假情试探。 因你也将真心真意瞒
        起来，我也将真心真意瞒起来，都只用假意试探，如此“两假相逢，终有一真”， 其间琐琐碎碎，难保不有口角之事。}
}
{
    \eA{And choose a few from among those young pages to remain on duty, at the
        second gate and at the two side entrances, so as to ask for things and
        deliver messages. Do you understand me, yes or no? The young ladies and
        ladies have all come out of town to-day, and not a single outsider must be
        permitted to put his foot in here." "I understand," replied Lin Chih-hsiao
        hurriedly signifying his obedience.}
}
{
}

\Verse{34}
{
    \zA{即如此刻，宝玉的心内想的是：“别人不知我 的心还可恕，难道你就不想我的心里眼里只有你？ 你不能为我解烦恼，反来拿这个
        话堵噎我，可见我心里时时刻刻白有你，你心里竟没我了。” 宝玉是这个意思，只 口里说不出来。 那黛玉心里想着：“你心里自然有我，虽有‘金玉相对’之说，你
        岂是重这邪说不重人的呢？ 我就时常提这‘金玉’，你只管了然无闻的，方见的是 待我重，无毫发私心了。 怎么我只一提‘金玉’的事，你就着急呢？}
}
{
    \eA{Next he uttered several yes's. "Now," proceeded Chia Chen; "you can go on
        your way. But how is it, I don't see anything of Jung Erh?" he went on to
        ask. This question was barely out of his lips, when he caught sight of Jung
        Erh running out of the belfry. "Look at him," shouted Chia Chen. "Look at
        him! I don't feel hot in here, and yet he must go in search of a cool place.
        Spit at him!"}
}
{
}

\Verse{35}
{
    \zA{可知你心里时 时有这个‘金玉’的念头。 我一提，你怕我多心，故意儿着急，安心哄我。” 那宝
        玉心中又想着：“我不管怎么样都好，只要你随意，我就立刻因你死了，也是情愿 的。 你知也罢，不知也罢，只由我的心，那才是你和我近，不和我远。” 黛玉心里
        又想着：“你只管你就是了。 你好，我自然好。 你要把自己丢开，只管周旋我，是 你不叫我近你，竟叫我远了。”
        看官，你道两个人原是一个心，如此看来，却都是多生了枝叶，将那求近之心 反弄成疏远之意了。}
}
{
    \eA{he cried to the family servants. The young pages were fully aware that Chia
        Chen's ordinary disposition was such that he could not brook contradiction,
        and one of the lads speedily came forward and sputtered in Chia Jung's face.}
}
{
}

\Verse{36}
{
    \zA{此皆他二人素昔所存私心，难以备述。 如今只说他们外面的形 容。 那宝玉又听见他说“好姻缘”三个字，越发逆了己意。 心里干噎，口里说不出
        来，便赌气向颈上摘下通灵玉来，咬咬牙，狠命往地下一摔，道：“什么劳什子! 我砸了你，就完了事了！” 偏生那玉坚硬非常，摔了一下，竟文风不动。 宝玉见不
        破，便回身找东西来砸。 黛玉见他如此，早已哭起来，说道：“何苦来你砸那哑吧 东西？ 有砸他的，不如来砸我！” 二人闹着，紫鹃雪雁等忙来解劝。}
}
{
    \eA{But Chia Chen still kept his gaze fixed on him, so the young page had to
        inquire of Chia Jung: "Master doesn't feel hot here, and how is it that you,
        Sir, have been the first to go and get cool?" Chia Jung however dropped his
        arms, and did not venture to utter a single sound.}
}
{
}

\Verse{37}
{
    \zA{后来见宝玉下死劲的砸那玉，忙上来夺，又 夺不下来。 见比往日闹的大了，少不得去叫袭人。 袭人忙赶了来，才夺下来。 宝玉
        冷笑道：“我是砸我的东西，与你们什么相干！” 袭人见他脸都气黄了，眉眼都变 了，从来没气的这么样，便拉着他的手，笑道：“你合妹妹拌嘴，不犯着砸他； 倘
        或砸坏了，叫他心里脸上怎么过的去呢？” 黛玉一行哭着，一行听了这话，说到自 己心坎儿上来，可见宝玉连袭人不如，越发伤心大哭起来。}
}
{
    \eA{Chia Yün, Chia P'ing, Chia Ch'in and the other young people overheard what
        was going on and not only were they scared out of their wits, but even Chia
        Lien, Chia Pin, Chia Ch'ung and their companions were stricken with intense
        fright and one by one they quietly slipped down along the foot of the wall.
        "What are you standing there for?" Chia Chen shouted to Chia Jung.}
}
{
}

\Verse{38}
{
    \zA{心里一急，方才吃的香 薷饮，便承受不住，“哇”的一声，都吐出来了。 紫鹃忙上来用绢子接住，登时一 口一口的，把块绢子吐湿。 雪雁忙上来捶揉。
        紫鹃道：“虽然生气，姑娘到底也该 保重些。 才吃了药，好些儿，这会子因和宝二爷拌嘴，又吐出来了； 倘或犯了病， 宝二爷怎么心里过的去呢？”
        宝玉听了这话，说到自己心坎儿上来，可见黛玉竟还 不如紫鹃呢。 又见黛玉脸红头胀，一行啼哭，一行气凑，一行是泪，一行是汗，不 胜怯弱。}
}
{
    \eA{"Don't you yet get on your horse and gallop home and tell your mother that
        our venerable senior is here with all the young ladies, and bid them come at
        once and wait upon them?" As soon as Chia Jung heard these words, he ran out
        with hurried stride and called out repeatedly for his horse.}
}
{
}

\Verse{39}
{
    \zA{宝玉见了这般，又自己后悔：“方才不该和他较证，这会子他这样光景， 我又替不了他。” 心里想着，也由不得滴下泪来了。
        袭人守着宝玉，见他两个哭的悲痛，也心酸起来。 又摸着宝玉的手冰凉，要劝 宝玉不哭罢，一则恐宝玉有什么委屈闷在心里，二则又恐薄了黛玉：两头儿为难。
        正是女儿家的心性，不觉也流下泪来。 紫鹃一面收拾了吐的药，一面拿扇子替黛玉 轻轻的扇着，见三个人都鸦雀无声，各自哭各自的，索性也伤起心来，也拿着绢子
        拭泪。 四个人都无言对泣。}
}
{
    \eA{Now he felt resentment, arguing within himself: "Who knows what he has been
        up to the whole morning, that he now finds fault with me!" Now he went on to
        abuse the young servants, crying: "Are your hands made fast, that you can't
        lead the horse round?"}
}
{
}

\Verse{40}
{
    \zA{还是袭人勉强笑向宝玉道：“你不看别的，你看看这玉 上穿的穗子，也不该和林姑娘拌嘴呀。” 黛玉听了，也不顾病，赶来夺过去，顺手 抓起一把剪子来就铰。
        袭人紫鹃刚要夺，已经剪了几段。 黛玉哭道：“我也是白效 力，他也不稀罕，自有别人替他再穿好的去呢！” 袭人忙接了玉道：“何苦来!这
        是我才多嘴的不是了。” 宝玉向黛玉道：“你只管铰!我横竖不带他，也没什么。”}
}
{
    \eA{And he felt inclined to bid a servant-boy go on the errand, but fearing
        again lest he should subsequently be found out, and be at a loss how to
        account for his conduct he felt compelled to proceed in person; so mounting
        his steed, he started on his way. But to return to Chia Chen. Just as he was
        about to be take himself inside, he noticed the Taoist Chang, who stood next
        to him, force a smile.}
}
{
}

\Verse{41}
{
    \zA{只顾里头闹，谁知那些老婆子们见黛玉大哭大吐，宝玉又砸玉，不知道要闹到 什么田地儿，便连忙的一齐往前头去回了贾母王夫人知道，好不至于连累了他们。
        那贾母王夫人见他们忙忙的做一件正经事来告诉，也都不知有了什么原故，便一齐 进园来瞧。 急的袭人抱怨紫鹃：“为什么惊动了老太太、太太？” 紫鹃又只当是袭
        人着人去告诉的，也抱怨袭人。}
}
{
    \eA{"I'm not properly speaking," he remarked, "on the same footing as the others
        and should be in attendance inside, but as on account of the intense heat,
        the young ladies have come out of doors, I couldn't presume to take upon
        myself to intrude and ask what your orders, Sir, are. But the dowager lady
        may possibly inquire about me, or may like to visit any part of the temple,
        so I shall wait in here."}
}
{
}

\Verse{42}
{
    \zA{那贾母王夫人进来，见宝玉也无言，黛玉也无话， 问起来，又没为什么事，便将这祸移到袭人紫鹃两个人身上，说：“为什么你们不
        小心伏侍，这会子闹起来都不管呢？” 因此将二人连骂带说教训了一顿。 二人都没 的说，只得听着。 还是贾母带出宝玉去了，方才平伏。
        过了一日，至初三日，乃是薛蟠生日，家里摆酒唱戏，贾府诸人都去了。 宝玉 因得罪了黛玉，二人总未见面，心中正自后悔，无精打彩，那里还有心肠去看戏，
        因而推病不去。}
}
{
    \eA{Chia Chen was fully cognisant that this Taoist priest, Chang, had, it is
        true, in past days, stood as a substitute for the Duke of the Jung Kuo
        mansion, but that the former Emperor had, with his own lips, conferred upon
        him the appellation of the 'Immortal being of the Great Unreal,' that he
        held at present the seal of 'Taoist Superior,' that the reigning Emperor had
        raised him to the rank of the 'Pure man,' that the princes, now-a-days,
        dukes, and high officials styled him the "Supernatural being," and he did
        not therefore venture to treat him with any disrespect.}
}
{
}

\Verse{43}
{
    \zA{黛玉不过前日中了些暑溽之气，本无甚大病，听见他不去，心里想： “他是好吃酒听戏的，今日反不去，自然是因为昨儿气着了； 再不然他见我不去， 他也没心肠去。
        只是昨儿千不该万不该铰了那玉上的穗子。 管定他再不带了，还得 我穿了他才带。” 因而心中十分后悔。 那贾母见他两个都生气，只说趁今儿那边去
        看戏，他两个见了，也就完了，不想又都不去。 老人家急的抱怨说：“我这老冤家， 是那一世里造下的孽障？}
}
{
    \eA{In the second place, (he knew that) he had paid frequent visits to the
        mansions, and that he had made the acquaintance of the ladies and young
        ladies, so when he heard his present remark he smilingly rejoined. "Do you
        again make use of such language amongst ourselves? One word more, and I'll
        take that beard of yours, and outroot it! Don't you yet come along with me
        inside?"}
}
{
}

\Verse{44}
{
    \zA{偏偏儿的遇见了这么两个不懂事的小冤家儿，没有一天不 叫我操心!真真的是俗语儿说的，‘不是冤家不聚头’了。 几时我闭了眼，断了这
        口气，任凭你们两个冤家闹上天去，我‘眼不见，心不烦’，也就罢了。 偏他娘的 又不咽这口气！” 自己抱怨着，也哭起来了。 谁知这个话传到宝玉黛玉二人耳内，
        他二人竟从来没有听见过“不是冤家不聚头”的这句俗话儿，如今忽然得了这句话， 好似参禅的一般，都低着头细嚼这句话的滋味儿，不觉的潸然泪下。}
}
{
    \eA{"Hah, hah," laughed the Taoist Chang aloud, as he followed Chia Chen in.
        Chia Chen approached dowager lady Chia. Bending his body he strained a
        laugh. "Grandfather Chang," he said, "has come in to pay his respects."
        "Raise him up!" old lady Chia vehemently called out. Chia Chen lost no time
        in pulling him to his feet and bringing him over. The Taoist Chang first
        indulged in loud laughter.}
}
{
}

\Verse{45}
{
    \zA{虽然不曾会面， 却一个在潇湘馆临风洒泪，一个在怡红院对月长吁，正是“人居两地，情发一心” 了。 袭人因劝宝玉道：“千万不是，都是你的不是。
        往日家里的小厮们和他的姐姐 妹妹拌嘴，或是两口子分争，你要是听见了，还骂那些小厮们蠢，不能体贴女孩儿 们的心肠； 今儿怎么你也这么着起来了？
        明儿初五，大节下的，你们两个再这么仇 人似的，老太太越发要生气了，一定弄的大家不安生。}
}
{
    \eA{"Oh Buddha of unlimited years!" he then observed. "Have you kept all right
        and in good health, throughout, venerable Senior? Have all the ladies and
        young ladies continued well? I haven't been for some time to your mansion to
        pay my obeisance, but you, my dowager lady, have improved more and more."
        "Venerable Immortal Being!" smiled old lady Chia, "how are you; quite well?"}
}
{
}

\Verse{46}
{
    \zA{依我劝你，正经下个气儿， 赔个不是，大家还是照常一样儿的，这么着不好吗？” 宝玉听了，不知依与不依。 要知端详，下回分解。 (本章完)}
}
{
    \eA{"Thanks to the ten thousand blessings he has enjoyed from your hands,"
        rejoined Chang the Taoist, "your servant too continues pretty strong and
        hale. In every other respect, I've, after all, been all right; but I have
        felt much concern about Mr. Pao-yü. Has he been all right all the time? The
        other day, on the 26th of the fourth moon, I celebrated the birthday of the
        'Heaven-Pervading-Mighty-King;'}
}
{
}

\Verse{47}
{
    \zA{}
}
{
    \eA{few people came and everything went off right and proper. I told them to
        invite Mr. Pao to come for a stroll; but how was it they said that he wasn't
        at home?" "It was indeed true that he was away from home," remarked dowager
        lady Chia. As she spoke, she turned her head round and called Pao-yü. Pao-yü
        had, as it happened, just returned from outside where he had been to make
        himself comfortable, and with speedy step, he came forward. "My respects to
        you, grandfather Chang," he said. The Taoist Chang eagerly clasped him in
        his arms and inquired how he was getting on. Turning towards old lady Chia,
        "Mr. Pao," he observed, "has grown fatter than ever."}
}
{
}

\Verse{48}
{
    \zA{}
}
{
    \eA{"Outwardly, his looks," replied dowager lady Chia, "may be all right, but,
        inwardly, he is weak. In addition to this, his father presses him so much to
        study that he has again and again managed, all through this bullying, to
        make his child fall sick." "The other day," continued Chang the Taoist, "I
        went to several places on a visit, and saw characters written by Mr. Pao and
        verses composed by him, all of which were exceedingly good; so how is it
        that his worthy father still feels displeased with him, and maintains that
        Mr. Pao is not very fond of his books? According to my humble idea, he knows
        quite enough.}
}
{
}

\Verse{49}
{
    \zA{}
}
{
    \eA{As I consider Mr. Pao's face, his bearing, his speech and his deportment,"
        he proceeded, heaving a sigh, "what a striking resemblance I find in him to
        the former duke of the Jung mansion!" As he uttered these words, tears
        rolled down his cheeks. At these words, old lady Chia herself found it hard
        to control her feelings. Her face became covered with the traces of tears.
        "Quite so," she assented, "I've had ever so many sons and grandsons, and not
        one of them betrayed the slightest resemblance to his grandfather; and this
        Pao-yü turns out to be the very image of him!"}
}
{
}

\Verse{50}
{
    \zA{}
}
{
    \eA{"What the former duke of Jung Kuo was like in appearance," Chang, the Taoist
        went on to remark, addressing himself to Chia Chen, "you gentlemen, and your
        generation, were, of course, needless to say, not in time to see for
        yourselves; but I fancy that even our Senior master and our Master Secundus
        have but a faint recollection of it." This said, he burst into another loud
        fit of laughter. "The other day," he resumed, "I was at some one's house and
        there I met a young girl, who is this year in her fifteenth year, and verily
        gifted with a beautiful face, and I bethought myself that Mr. Pao must also
        have a wife found for him. As far as looks, intelligence and mental talents,
        extraction and family standing go, this maiden is a suitable match for him.}
}
{
}

\Verse{51}
{
    \zA{}
}
{
    \eA{But as I didn't know what your venerable ladyship would have to say about
        it, your servant did not presume to act recklessly, but waited until I could
        ascertain your wishes before I took upon myself to open my mouth with the
        parties concerned." "Some time ago," responded dowager lady Chia, "a bonze
        explained that it was ordained by destiny that this child shouldn't be
        married at an early age, and that we should put things off until he grew
        somewhat in years before anything was settled. But mark my words now.}
}
{
}

\Verse{52}
{
    \zA{}
}
{
    \eA{Pay no regard as to whether she be of wealthy and honourable stock or not,
        the essential thing is to find one whose looks make her a fit match for him
        and then come at once and tell me. For even admitting that the girl is poor,
        all I shall have to do will be to bestow on her a few ounces of silver; but
        fine looks and a sweet temperament are not easy things to come across." When
        she had done speaking, lady Feng was heard to smilingly interpose:
        "Grandfather Chang, aren't you going to change the talisman of 'Recorded
        Name' of our daughter? The other day, lucky enough for you, you had again
        the great cheek to send some one to ask me for some satin of gosling-yellow
        colour. I gave it to you, for had I not, I was afraid lest your old face
        should have been made to feel uneasy."}
}
{
}

\Verse{53}
{
    \zA{}
}
{
    \eA{"Hah, hah," roared the Taoist Chang, "just see how my eyes must have grown
        dim! I didn't notice that you, my lady, were in here; nor did I express one
        word of thanks to you! The talisman of 'Recorded Name' is ready long ago. I
        meant to have sent it over the day before yesterday, but the unforeseen
        visit of the Empress to perform meritorious deeds upset my equilibrium, and
        made me quite forget it. But it's still placed before the gods, and if you
        will wait I'll go and fetch it." Saying this, he rushed into the main hall.
        Presently, he returned with a tea-tray in hand, on which was spread a deep
        red satin cover, brocaded with dragons. In this, he presented the charm. Ta
        Chieh-erh's nurse took it from him.}
}
{
}

\Verse{54}
{
    \zA{}
}
{
    \eA{But just as the Taoist was on the point of taking Ta Chieh-erh in his
        embrace, lady Feng remarked with a smile: "It would have been sufficient if
        you'd carried it in your hand! And why use a tray to lay it on?" "My hands
        aren't clean," replied the Taoist Chang, "so how could I very well have
        taken hold of it? A tray therefore made things much cleaner!" "When you
        produced that tray just now," laughed lady Feng, "you gave me quite a start;
        I didn't imagine that it was for the purpose of bringing the charm in. It
        really looked as if you were disposed to beg donations of us." This
        observation sent the whole company into a violent fit of laughter. Even Chia
        Chen could not suppress a smile. "What a monkey!" dowager lady Chia
        exclaimed, turning her head round.}
}
{
}

\Verse{55}
{
    \zA{}
}
{
    \eA{"What a monkey you are! Aren't you afraid of going down to that Hell, where
        tongues are cut off?" "I've got nothing to do with any men whatever,"
        rejoined lady Feng laughing, "and why does he time and again tell me that
        it's my bounden duty to lay up a store of meritorious deeds; and that if I'm
        remiss, my life will be short?" Chang, the Taoist, indulged in further
        laughter. "I brought out," he explained, "the tray so as to kill two birds
        with one stone. It wasn't, however, to beg for donations. On the contrary,
        it was in order to put in it the jade, which I meant to ask Mr. Pao to take
        off, so as to carry it outside and let all those Taoist friends of mine, who
        come from far away, as well as my neophytes and the young apprentices, see
        what it's like."}
}
{
}

\Verse{56}
{
    \zA{}
}
{
    \eA{"Well, since that be the case," added old lady Chia, "why do you, at your
        age, try your strength by running about the whole day long? Take him at once
        along and let them see it! But were you to have called him in there,
        wouldn't it have saved a lot of trouble?" "Your venerable ladyship," resumed
        Chang, the Taoist, "isn't aware that though I be, to look at, a man of
        eighty, I, after all, continue, thanks to your protection, my dowager lady,
        quite hale and strong. In the second place, there are crowds of people in
        the outer rooms; and the smells are not agreeable. Besides it's a very hot
        day and Mr. Pao couldn't stand the heat as he is not accustomed to it.}
}
{
}

\Verse{57}
{
    \zA{}
}
{
    \eA{So were he to catch any disease from the filthy odours, it would be a grave
        thing!" After these forebodings old lady Chia accordingly desired Pao-yü to
        unclasp the jade of Spiritual Perception, and to deposit it in the tray. The
        Taoist, Chang, carefully ensconced it in the folds of the wrapper,
        embroidered with dragons, and left the room, supporting the tray with both
        his hands. During this while, dowager lady Chia and the other inmates
        devoted more of their time in visiting the various places. But just as they
        were on the point of going up the two-storied building, they heard Chia Chen
        shout: "Grandfather Chang has brought back the jade." As he spoke, the
        Taoist Chang was seen advancing up to them, the tray in hand. "The whole
        company," he smiled, "were much obliged to me.}
}
{
}

\Verse{58}
{
    \zA{}
}
{
    \eA{They think Mr. Pao's jade really lovely! None of them have, however, any
        suitable gifts to bestow. These are religious articles, used by each of them
        in propagating the doctrines of Reason, but they're all only too ready to
        give them as congratulatory presents. If, Mr. Pao, you don't fancy them for
        anything else, just keep them to play with or to give away to others."
        Dowager lady Chia, at these words, looked into the tray. She discovered that
        its contents consisted of gold signets, and jade rings, or sceptres,
        implying: "may you have your wishes accomplished in everything," or "may you
        enjoy peace and health from year to year;" that the various articles were
        strung with pearls or inlaid with precious stones, worked in jade or mounted
        in gold;}
}
{
}

\Verse{59}
{
    \zA{}
}
{
    \eA{and that they were in all from thirty to fifty. "What nonsense you're
        talking!" she then exclaimed. "Those people are all divines, and where could
        they have rummaged up these things? But what need is there for any such
        presents? He may, on no account, accept them." "These are intended as a
        small token of their esteem," responded Chang, the Taoist, smiling, "your
        servant cannot therefore venture to interfere with them. If your venerable
        ladyship will not keep them, won't you make it patent to them that I'm
        treated contemptuously, and unlike what one should be, who has joined the
        order through your household?"}
}
{
}

\Verse{60}
{
    \zA{}
}
{
    \eA{Only when old lady Chia heard these arguments did she direct a servant to
        receive the presents. "Venerable senior," Pao-yü smilingly chimed in. "After
        the reasons advanced by grandfather Chang, we cannot possibly refuse them.
        But albeit I feel disposed to keep these things, they are of no avail to me;
        so would it not be well were a servant told to carry the tray and to follow
        me out of doors, that I may distribute them to the poor? "You are perfectly
        right in what you say!" smiled dowager lady Chia. The Taoist Chang, however,
        went on speedily to use various arguments to dissuade him. "Mr. Pao," he
        observed, "your intention is, it is true, to perform charitable acts;}
}
{
}

\Verse{61}
{
    \zA{}
}
{
    \eA{but though you may aver that these things are of little value, you'll
        nevertheless find among them several articles you might turn to some
        account. Were you to let the beggars have them, why they will, first of all,
        be none the better for them; and, next, it will contrariwise be tantamount
        to throwing them away! If you want to distribute anything among the poor,
        why don't you dole out cash to them?" "Put them by!" promptly shouted
        Pao-yü, after this rejoinder, "and when evening comes, take a few cash and
        distribute them." These directions given, Chang, the Taoist, retired out of
        the place. Dowager lady Chia and her companions thereupon walked upstairs
        and sat in the main part of the building.}
}
{
}

\Verse{62}
{
    \zA{}
}
{
    \eA{Lady Feng and her friends adjourned into the eastern part, while the
        waiting-maids and servants remained in the western portion, and took their
        turns in waiting on their mistresses. Before long, Chia Chen came back. "The
        plays," he announced, "have been chosen by means of slips picked out before
        the god. The first one on the list is the 'Record of the White Snake.'" "Of
        what kind of old story does 'the record of the white snake,' treat?" old
        lady Chia inquired. "The story about Han Kao-tsu," replied Chia Chen,
        "killing a snake and then ascending the throne. The second play is, 'the Bed
        covered with ivory tablets.'" "Has this been assigned the second place?"
        asked dowager lady Chia. "Yet never mind; for as the gods will it thus,
        there is no help than not to demur.}
}
{
}

\Verse{63}
{
    \zA{}
}
{
    \eA{But what about the third play?" she went on to inquire. "The Nan Ko dream is
        the third," Chia Chen answered. This response elicited no comment from
        dowager lady Chia. Chia Chen therefore withdrew downstairs, and betook
        himself outside to make arrangements for the offerings to the gods, for the
        paper money and eatables that had to be burnt, and for the theatricals about
        to begin. So we will leave him without any further allusion, and take up our
        narrative with Pao-yü. Seating himself upstairs next to old lady Chia, he
        called to a servant-girl to fetch the tray of presents given to him a short
        while back, and putting on his own trinket of jade, he fumbled about with
        the things for a bit, and picking up one by one, he handed them to his
        grandmother to admire.}
}
{
}

\Verse{64}
{
    \zA{}
}
{
    \eA{But old lady Chia espied among them a unicorn, made of purplish gold, with
        kingfisher feathers inserted, and eagerly extending her arm, she took it up.
        "This object," she smiled, "seems to me to resemble very much one I've seen
        worn also by the young lady of some household or other of ours." "Senior
        cousin, Shih Hsiang-yün," chimed in Pao-ch'ai, a smile playing on her lips,
        "has one, but it's a trifle smaller than this." "Is it indeed Yün-erh who
        has it?" exclaimed old lady Chia. "Now that she lives in our house,"
        remarked Pao-yü, "how is it that even I haven't seen anything of it?"
        "Cousin Pao-ch'ai," rejoined T'an Ch'un laughingly, "has the power of
        observation; no matter what she sees, she remembers."}
}
{
}

\Verse{65}
{
    \zA{}
}
{
    \eA{Lin Tai-yü gave a sardonic smile. "As far as other matters are concerned,"
        she insinuated, "her observation isn't worth speaking of; where she's
        extra-observant is in articles people may wear about their persons."
        Pao-chai, upon catching this sneering remark, at once turned her head round,
        and pretended she had not heard. But as soon as Pao-yü learnt that Shih
        Hsiang-yün possessed a similar trinket, he speedily picked up the unicorn,
        and hid it in his breast, indulging, at the same time, in further
        reflection. Yet, fearing lest people might have noticed that he kept back
        that particular thing the moment he discovered that Shih Hsiang-yün had one
        identical with it, he fixed his eyes intently upon all around while
        clutching it.}
}
{
}

\Verse{66}
{
    \zA{}
}
{
    \eA{He found however that not one of them was paying any heed to his movements
        except Lin Tai-yü, who, while gazing at him was, nodding her head, as if
        with the idea of expressing her admiration. Pao-yü, therefore, at once felt
        inwardly ill at ease, and pulling out his hand, he observed, addressing
        himself to Tai-yü with an assumed smile, "This is really a fine thing to
        play with; I'll keep it for you, and when we get back home, I'll pass a
        ribbon through it for you to wear." "I don't care about it," said Lin
        Tai-yü, giving her head a sudden twist. "Well," continued Pao-yü laughingly,
        "if you don't like it, I can't do otherwise than keep it myself." Saying
        this, he once again thrust it away.}
}
{
}

\Verse{67}
{
    \zA{}
}
{
    \eA{But just as he was about to open his lips to make some other observation, he
        saw Mrs. Yu, the spouse of Chia Chen, arrive along with the second wife
        recently married by Chia Jung, that is, his mother and her daughter-in-law,
        to pay their obeisance to dowager lady Chia. "What do you people rush over
        here for again?" old lady Chia inquired. "I came here for a turn, simply
        because I had nothing to do." But no sooner was this inquiry concluded than
        they heard a messenger announce: "that some one had come from the house of
        general Feng." The family of Feng Tzu-ying had, it must be explained, come
        to learn the news that the inmates of the Chia mansion were offering a
        thanksgiving service in the temple, and, without loss of time, they got
        together presents of pigs, sheep, candles, tea and eatables and sent them
        over.}
}
{
}

\Verse{68}
{
    \zA{}
}
{
    \eA{The moment lady Feng heard about it she hastily crossed to the main part of
        the two-storied building. "Ai-ya;" she ejaculated, clapping her hands and
        laughing. "I never expected anything of the sort; we merely said that we
        ladies were coming for a leisurely stroll and people imagined that we were
        spreading a sumptuous altar with lenten viands and came to bring us
        offerings! But it's all our old lady's fault for bruiting it about! Why, we
        haven't even got any slips of paper with tips ready." She had just finished
        speaking, when she perceived two matrons, who acted as house-keepers in the
        Feng family, walk upstairs. But before the Feng servants could take their
        leave, presents likewise arrived, in quick succession, from Chao, the
        Vice-President of the Board.}
}
{
}

\Verse{69}
{
    \zA{}
}
{
    \eA{In due course, one lot of visitors followed another. For as every one got
        wind of the fact that the Chia family was having thanksgiving services, and
        that the ladies were in the temple, distant and close relatives, friends,
        old friends and acquaintances all came to present their contributions. So
        much so, that dowager lady Chia began at this juncture to feel sorry that
        she had ever let the cat out of the bag. "This is no regular fasting," she
        said, "we simply have come for a little change; and we should not have put
        any one to any inconvenience!" Although therefore she was to have remained
        present all day at the theatrical performance, she promptly returned home
        soon after noon, and the next day she felt very loth to go out of doors
        again. "By striking the wall, we've also stirred up dust," lady Feng argued.}
}
{
}

\Verse{70}
{
    \zA{}
}
{
    \eA{"Why we've already put those people to the trouble so we should only be too
        glad to-day to have another outing." But as when dowager lady Chia
        interviewed the Taoist Chang, the previous day, he made allusion to Pao-yü
        and canvassed his engagement, Pao-yü experienced, little as one would have
        thought it, much secret displeasure during the whole of that day, and on his
        return home he flew into a rage and abused Chang, the rationalistic priest,
        for harbouring designs to try and settle a match for him.}
}
{
}

\Verse{71}
{
    \zA{}
}
{
    \eA{At every breath and at every word he resolved that henceforward he would not
        set eyes again upon the Taoist Chang. But no one but himself had any idea of
        the reason that actuated him to absent himself. In the next place, Lin
        Tai-yü began also, on her return the day before, to ail from a touch of the
        sun, so their grandmother was induced by these two considerations to remain
        firm in her decision not to go. When lady Feng, however, found that she
        would not join them, she herself took charge of the family party and set out
        on the excursion. But without descending to particulars, let us advert to
        Pao-yü. Seeing that Lin Tai-yü had fallen ill, he was so full of solicitude
        on her account that he even had little thought for any of his meals, and not
        long elapsed before he came to inquire how she was.}
}
{
}

\Verse{72}
{
    \zA{}
}
{
    \eA{Tai-yü, on her part, gave way to fear lest anything should happen to him,
        (and she tried to re-assure him). "Just go and look at the plays," she
        therefore replied, "what's the use of boxing yourself up at home?" Pao-yü
        was, however, not in a very happy frame of mind on account of the reference
        to his marriage made by Chang, the Taoist, the day before, so when he heard
        Lin Tai-yü's utterances: "If others don't understand me;" he mused, "it's
        anyhow excusable; but has she too begun to make fun of me?"}
}
{
}

\Verse{73}
{
    \zA{}
}
{
    \eA{His heart smarted in consequence under the sting of a mortification a
        hundred times keener than he had experienced up to that occasion. Had he
        been with any one else, it would have been utterly impossible for her to
        have brought into play feelings of such resentment, but as it was no other
        than Tai-yü who spoke the words, the impression produced upon him was indeed
        different from that left in days gone by, when others employed similar
        language. Unable to curb his feelings, he instantaneously lowered his face.}
}
{
}

\Verse{74}
{
    \zA{}
}
{
    \eA{"My friendship with you has been of no avail" he rejoined. "But, never mind,
        patience!" This insinuation induced Lin Tai-yü to smile a couple of
        sarcastic smiles. "Yes, your friendship with me has been of no avail," she
        repeated; "for how can I compare with those whose manifold qualities make
        them fit matches for you?" As soon as this sneer fell on Pao-yü's ear he
        drew near to her. "Are you by telling me this," he asked straight to her
        face, "deliberately bent upon invoking imprecations upon me that I should be
        annihilated by heaven and extinguished by earth?" Lin Tai-yü could not for a
        time fathom the import of his remarks. "It was," Pao-yü then resumed, "on
        account of this very conversation that I yesterday swore several oaths, and
        now would you really make me repeat another one?}
}
{
}

\Verse{75}
{
    \zA{}
}
{
    \eA{But were the heavens to annihilate me and the earth to extinguish me, what
        benefit would you derive?" This rejoinder reminded Tai-yü of the drift of
        their conversation on the previous day. And as indeed she had on this
        occasion framed in words those sentiments, which should not have dropped
        from her lips, she experienced both annoyance and shame, and she tremulously
        observed: "If I entertain any deliberate intention to bring any harm upon
        you, may I too be destroyed by heaven and exterminated by earth! But what's
        the use of all this!}
}
{
}

\Verse{76}
{
    \zA{}
}
{
    \eA{I know very well that the allusion to marriage made yesterday by Chang, the
        Taoist, fills you with dread lest he might interfere with your choice. You
        are inwardly so irate that you come and treat me as your malignant
        influence." Pao-yü, the fact is, had ever since his youth developed a
        peculiar kind of mean and silly propensity. Having moreover from tender
        infancy grown up side by side with Tai-Yü, their hearts and their feelings
        were in perfect harmony. More, he had recently come to know to a great
        extent what was what, and had also filled his head with the contents of a
        number of corrupt books and licentious stories.}
}
{
}

\Verse{77}
{
    \zA{}
}
{
    \eA{Of all the eminent and beautiful girls that he had met too in the families
        of either distant or close relatives or of friends, not one could reach the
        standard of Lin Tai-yü. Hence it was that he commenced, from an early period
        of his life, to foster sentiments of love for her; but as he could not very
        well give utterance to them, he felt time and again sometimes elated,
        sometimes vexed, and wont to exhaust every means to secretly subject her
        heart to a test. Lin Tai-yü happened, on the other hand, to possess in like
        manner a somewhat silly disposition; and she too frequently had recourse to
        feigned sentiments to feel her way.}
}
{
}

\Verse{78}
{
    \zA{}
}
{
    \eA{And as she began to conceal her true feelings and inclinations and to simply
        dissimulate, and he to conceal his true sentiments and wishes and to
        dissemble, the two unrealities thus blending together constituted eventually
        one reality. But it was hardly to be expected that trifles would not be the
        cause of tiffs between them. Thus it was that in Pao-yü's mind at this time
        prevailed the reflection: "that were others unable to read my feelings, it
        would anyhow be excusable; but is it likely that you cannot realise that in
        my heart and in my eyes there is no one else besides yourself.}
}
{
}

\Verse{79}
{
    \zA{}
}
{
    \eA{But as you were not able to do anything to dispel my annoyance, but made
        use, instead, of the language you did to laugh at me, and to gag my mouth,
        it's evident that though you hold, at every second and at every moment, a
        place in my heart, I don't, in fact, occupy a place in yours." Such was the
        construction attached to her conduct by Pao-yü, yet he did not have the
        courage to tax her with it. "If, really, I hold a place in your heart," Lin
        Tai-yü again reflected, "why do you, albeit what's said about gold and jade
        being a fit match, attach more importance to this perverse report and think
        nothing of what I say?}
}
{
}

\Verse{80}
{
    \zA{}
}
{
    \eA{Did you, when I so often broach the subject of this gold and jade, behave as
        if you, verily, had never heard anything about it, I would then have seen
        that you treat me with preference and that you don't harbour the least
        particle of a secret design. But how is it that the moment I allude to the
        topic of gold and jade, you at once lose all patience? This is proof enough
        that you are continuously pondering over that gold and jade, and that as
        soon as you hear me speak to you about them, you apprehend that I shall once
        more give way to conjectures, and intentionally pretend to be quite out of
        temper, with the deliberate idea of cajoling me!"}
}
{
}

\Verse{81}
{
    \zA{}
}
{
    \eA{These two cousins had, to all appearances, once been of one and the same
        mind, but the many issues, which had sprung up between them, brought about a
        contrary result and made them of two distinct minds. "I don't care what you
        do, everything is well," Pao-yü further argued, "so long as you act up to
        your feelings; and if you do, I shall be ever only too willing to even
        suffer immediate death for your sake. Whether you know this or not, doesn't
        matter; it's all the same. Yet were you to just do as my heart would have
        you, you'll afford me a clear proof that you and I are united by close ties
        and that you are no stranger to me!" "Just you mind your own business," Lin
        Tai-yü on her side cogitated. "If you will treat me well, I'll treat you
        well. And what need is there to put an end to yourself for my sake?}
}
{
}

\Verse{82}
{
    \zA{}
}
{
    \eA{Are you not aware that if you kill yourself, I'll also kill myself? But this
        demonstrates that you don't wish me to be near to you, and that you really
        want that I should be distant to you." It will thus be seen that the desire,
        by which they were both actuated, to strive and draw each other close and
        ever closer became contrariwise transformed into a wish to become more
        distant. But as it is no easy task to frame into words the manifold secret
        thoughts entertained by either, we will now confine ourselves to a
        consideration of their external manner. The three words "a fine match,"
        which Pao-yü heard again Lin Tai-yü pronounce proved so revolting to him
        that his heart got full of disgust and he was unable to give utterance to a
        single syllable.}
}
{
}

\Verse{83}
{
    \zA{}
}
{
    \eA{Losing all control over his temper, he snatched from his neck the jade of
        Spiritual Perception and, clenching his teeth, he spitefully dashed it down
        on the floor. "What rubbishy trash!" he cried. "I'll smash you to atoms and
        put an end to the whole question!" The jade, however, happened to be of
        extraordinary hardness, and did not, after all, sustain the slightest injury
        from this single fall. When Pao-yü realised that it had not broken, he
        forthwith turned himself round to get the trinket with the idea of carrying
        out his design of smashing it, but Tai-yü divined his intention, and soon
        started crying. "What's the use of all this!" she demurred, "and why, pray,
        do you batter that dumb thing about? Instead of smashing it, wouldn't it be
        better for you to come and smash me!"}
}
{
}

\Verse{84}
{
    \zA{}
}
{
    \eA{But in the middle of their dispute, Tzu Chüan, Hsüeh Yen and the other maids
        promptly interfered and quieted them. Subsequently, however, they saw how
        deliberately bent Pao-yü was upon breaking the jade, and they vehemently
        rushed up to him to snatch it from his hands. But they failed in their
        endeavours, and perceiving that he was getting more troublesome than he had
        ever been before, they had no alternative but to go and call Hsi Jen. Hsi
        Jen lost no time in running over and succeeded, at length, in getting hold
        of the trinket. "I'm smashing what belongs to me," remarked Pao-yü with a
        cynical smile, "and what has that to do with you people?"}
}
{
}

\Verse{85}
{
    \zA{}
}
{
    \eA{Hsi Jen noticed that his face had grown quite sallow from anger, that his
        eyes had assumed a totally unusual expression, and that he had never
        hitherto had such a fit of ill-temper and she hastened to take his hand in
        hers and to smilingly expostulate with him. "If you've had a tiff with your
        cousin," she said, "it isn't worth while flinging this down! Had you broken
        it, how would her heart and face have been able to bear the mortification?"
        Lin Tai-yü shed tears and listened the while to her remonstrances. Yet these
        words, which so corresponded with her own feelings, made it clear to her
        that Pao-yü could not even compare with Hsi Jen and wounded her heart so
        much more to the quick that she began to weep aloud.}
}
{
}

\Verse{86}
{
    \zA{}
}
{
    \eA{But the moment she got so vexed she found it hard to keep down the potion of
        boletus and the decoction, for counter-acting the effects of the sun, she
        had taken only a few minutes back, and with a retch she brought everything
        up. Tzu Chüan immediately pressed to her side and used her handkerchief to
        stop her mouth with. But mouthful succeeded mouthful, and in no time the
        handkerchief was soaked through and through. Hsüeh Yen then approached in a
        hurry and tapped her on the back. "You may, of course, give way to
        displeasure," Tzu Chüan argued; "but you should, after all, take good care
        of yourself Miss. You had just taken the medicines and felt the better for
        them; and here you now begin vomitting again; and all because you've had a
        few words with our master Secundus.}
}
{
}

\Verse{87}
{
    \zA{}
}
{
    \eA{But should your complaint break out afresh how will Mr. Pao bear the blow?"
        The moment Pao-yü caught this advice, which accorded so thoroughly with his
        own ideas, he found how little Tai-yü could hold her own with Tzu Chüan. And
        perceiving how flushed Tai-yü's face was, how her temples were swollen, how,
        while sobbing, she panted; and how, while crying, she was suffused with
        perspiration, and betrayed signs of extreme weakness, he began, at the sight
        of her condition, to reproach himself. "I shouldn't," he reflected, "have
        bandied words with her; for now that she's got into this frame of mind, I
        mayn't even suffer in her stead!"}
}
{
}

\Verse{88}
{
    \zA{}
}
{
    \eA{The self-reproaches, however, which gnawed his heart made it impossible for
        him to refrain from tears, much as he fought against them. Hsi Jen saw them
        both crying, and while attending to Pao-yü, she too unavoidably experienced
        much soreness of heart. She nevertheless went on rubbing Pao-yü's hands,
        which were icy cold. She felt inclined to advise Pao-yü not to weep, but
        fearing again lest, in the first place, Pao-yü might be inwardly aggrieved,
        and nervous, in the next, lest she should not be dealing rightly by Tai-yü,
        she thought it advisable that they should all have a good cry, as they might
        then be able to leave off. She herself therefore also melted into tears.}
}
{
}

\Verse{89}
{
    \zA{}
}
{
    \eA{As for Tzu-Chüan, at one time, she cleaned the expectorated medicine; at
        another, she took up a fan and gently fanned Tai-yü. But at the sight of the
        trio plunged in perfect silence, and of one and all sobbing for reasons of
        their own, grief, much though she did to struggle against it, mastered her
        feelings too, and producing a handkerchief, she dried the tears that came to
        her eyes. So there stood four inmates, face to face, uttering not a word and
        indulging in weeping. Shortly, Hsi Jen made a supreme effort, and smilingly
        said to Pao-yü: "If you don't care for anything else, you should at least
        have shown some regard for those tassels, strung on the jade, and not have
        wrangled with Miss Lin."}
}
{
}

\Verse{90}
{
    \zA{}
}
{
    \eA{Tai-yü heard these words, and, mindless of her indisposition, she rushed
        over, and snatching the trinket, she picked up a pair of scissors, lying
        close at hand, bent upon cutting the tassels. Hsi Jen and Tzu Chüan were on
        the point of wresting it from her, but she had already managed to mangle
        them into several pieces. "I have," sobbed Tai-yü, "wasted my energies on
        them for nothing; for he doesn't prize them. He's certain to find others to
        string some more fine tassels for him." Hsi Jen promptly took the jade. "Is
        it worth while going on in this way!" she cried. "But this is all my fault
        for having blabbered just now what should have been left unsaid." "Cut it,
        if you like!" chimed in Pao-yü, addressing himself to Tai-yü. "I will on no
        account wear it, so it doesn't matter a rap."}
}
{
}

\Verse{91}
{
    \zA{}
}
{
    \eA{But while all they minded inside was to create this commotion, they little
        dreamt that the old matrons had descried Tai-yü weep bitterly and vomit
        copiously, and Pao-yü again dash his jade on the ground, and that not
        knowing how far the excitement might not go, and whether they themselves
        might not become involved, they had repaired in a body to the front, and
        reported the occurrence to dowager lady Chia and Madame Wang, their object
        being to try and avoid being themselves implicated in the matter. Their old
        mistress and Madame Wang, seeing them make so much of the occurrence as to
        rush with precipitate haste to bring it to their notice, could not in the
        least imagine what great disaster might not have befallen them, and without
        loss of time they betook themselves together into the garden and came to see
        what the two cousins were up to.}
}
{
}

\Verse{92}
{
    \zA{}
}
{
    \eA{Hsi Jen felt irritated and harboured resentment against Tzu Chüan, unable to
        conceive what business she had to go and disturb their old mistress and
        Madame Wang. But Tzu Chüan, on the other hand, presumed that it was Hsi Jen,
        who had gone and reported the matter to them, and she too cherished angry
        feelings towards Hsi Jen. Dowager lady Chia and Madame Wang walked into the
        apartment. They found Pao-yü on one side saying not a word. Lin Tai-yü on
        the other uttering not a sound. "What's up again?" they asked. But throwing
        the whole blame upon the shoulders of Hsi Jen and Tzu Chüan, "why is it,"
        they inquired, "that you were not diligent in your attendance on them. They
        now start a quarrel, and don't you exert yourselves in the least to restrain
        them?"}
}
{
}

\Verse{93}
{
    \zA{}
}
{
    \eA{Therefore with obloquy and hard words they rated the two girls for a time in
        such a way that neither of them could put in a word by way of reply, but
        felt compelled to listen patiently. And it was only after dowager lady Chia
        had taken Pao-yü away with her that things quieted down again. One day
        passed. Then came the third of the moon. This was Hsüeh Pan's birthday, so
        in their house a banquet was spread and preparations made for a performance;
        and to these the various inmates of the Chia mansion went. But as Pao-yü had
        so hurt Tai-yü's feelings, the two cousins saw nothing whatever of each
        other, and conscience-stricken, despondent and unhappy, as he was at this
        time could he have had any inclination to be present at the plays?}
}
{
}

\Verse{94}
{
    \zA{}
}
{
    \eA{Hence it was that he refused to go on the pretext of indisposition. Lin
        Tai-yü had got, a couple of days back, but a slight touch of the sun and
        naturally there was nothing much the matter with her. When the news however
        reached her that he did not intend to join the party, "If with his weakness
        for wine and for theatricals," she pondered within herself, "he now chooses
        to stay away, instead of going, why, that quarrel with me yesterday must be
        at the bottom of it all. If this isn't the reason, well then it must be that
        he has no wish to attend, as he sees that I'm not going either.}
}
{
}

\Verse{95}
{
    \zA{}
}
{
    \eA{But I should on no account have cut the tassels from that jade, for I feel
        sure he won't wear it again. I shall therefore have to string some more on
        to it, before he puts it on." On this account the keenest remorse gnawed her
        heart. Dowager lady Chia saw well enough that they were both under the
        influence of temper. "We should avail ourselves of this occasion," she said
        to herself, "to go over and look at the plays, and as soon as the two young
        people come face to face, everything will be squared." Contrary to her
        expectations neither of them would volunteer to go. This so exasperated
        their old grandmother that she felt vexed with them.}
}
{
}

\Verse{96}
{
    \zA{}
}
{
    \eA{"In what part of my previous existence could an old sufferer like myself,"
        she exclaimed, "have incurred such retribution that my destiny is to come
        across these two troublesome new-fledged foes! Why, not a single day goes by
        without their being instrumental in worrying my mind! The proverb is indeed
        correct which says: 'that people who are not enemies are not brought
        together!' But shortly my eyes shall be closed, this breath of mine shall be
        snapped, and those two enemies will be free to cause trouble even up to the
        very skies; for as my eyes will then loose their power of vision, and my
        heart will be void of concern, it will really be nothing to me. But I
        couldn't very well stifle this breath of life of mine!" While inwardly a
        prey to resentment, she also melted into tears.}
}
{
}

\Verse{97}
{
    \zA{}
}
{
    \eA{These words were brought to the ears of Pao-yü and Tai-yü. Neither of them
        had hitherto heard the adage: "people who are not enemies are not brought
        together," so when they suddenly got to know the line, it seemed as if they
        had apprehended abstraction. Both lowered their heads and meditated on the
        subtle sense of the saying. But unconsciously a stream of tears rolled down
        their cheeks. They could not, it is true, get a glimpse of each other; yet
        as the one was in the Hsiao Hsiang lodge, standing in the breeze, bedewed
        with tears, and the other in the I Hung court, facing the moon and heaving
        deep sighs, was it not, in fact, a case of two persons living in two
        distinct places, yet with feelings emanating from one and the same heart?
        Hsi Jen consequently tendered advice to Pao-yü.}
}
{
}

\Verse{98}
{
    \zA{}
}
{
    \eA{"You're a million times to blame," she said, "it's you who are entirely at
        fault! For when some time ago the pages in the establishment, wrangled with
        their sisters, or when husband and wife fell out, and you came to hear
        anything about it, you blew up the lads, and called them fools for not
        having the heart to show some regard to girls; and now here you go and
        follow their lead. But to-morrow is the fifth day of the moon, a great
        festival, and will you two still continue like this, as if you were very
        enemies? If so, our venerable mistress will be the more angry, and she
        certainly will be driven sick! I advise you therefore to do what's right by
        suppressing your spite and confessing your fault, so that we should all be
        on the same terms as hitherto.}
}
{
}

\Verse{99}
{
    \zA{}
}
{
    \eA{You here will then be all right, and so will she over there." Pao-yü
        listened to what she had to say; but whether he fell in with her views or
        not is not yet ascertained; yet if you, reader, choose to know, we will
        explain in the next chapter.}
}
{
}
